% figures/w03-inner-loop.svg — 내부 루프 둘과 그 출력이 축 회전수가 되기까지.
%
%   bash _tools/tikz2svg.sh figures/src/w03-inner-loop.tex figures/w03-inner-loop.svg
%
% 사용자 수업 노트(2026-10-13) 1쪽 위의 구조를 그대로 옮긴 것이다. 사용자 지시:
% "inner loop 제어기에서 W04, W03 모두 강의 자료에 넣어줘... 내부 제어기가 속도제어,
%  헤딩제어이고, 출력이 X, N 이고 그것을 어떻게 control allocation 해서 nL, nR 로
%  분배하는지 공식까지 같이 보여주고... 각각 기능들이 잘 보이도록"
%
% 이 그림은 **3·4·5 주차가 함께 쓴다.** 그래서 한 장만 그린다.
%   3 주차  위 줄(속도 루프)이 주제. 아래 줄은 4 주차가 채운다
%   4 주차  아래 줄(선수각 루프)이 주제. ssa 와 요각속도 되먹임이 여기 있다
%   5 주차  왼쪽의 점선 상자가 채워진다 — figures/w05-guidance-loop.svg 가 그것이다
%
% 상자 안에 **그 블록이 푸는 식과 Simulink 블록 이름**을 함께 적는다. 구조만 그리면
% 학생이 도면과 모델을 잇지 못하고, 식만 적으면 어디에 사는 식인지 모른다 (§13).
%
% 배치 주의 (렌더해서 고친 것, 2026-10-13):
%   `minimum width` 는 **글이 길면 무시된다** — 상자가 글 폭만큼 자라 옆 상자를 덮었다.
%   폭을 정해야 하는 상자에는 반드시 `text width` 를 쓴다.
\documentclass[border=6pt]{standalone}
% gnc-style.tex — 이 볼트의 TikZ 그림이 공유하는 스타일.
%
% 저널 그림 규약을 스타일 하나로 굳혀 둔다. 그림마다 선 굵기나 화살촉을
% 다시 정하지 않는다 — 그것이 그림이 제각각으로 보이는 원인이다.
%
% 쓰는 법:  % gnc-style.tex — 이 볼트의 TikZ 그림이 공유하는 스타일.
%
% 저널 그림 규약을 스타일 하나로 굳혀 둔다. 그림마다 선 굵기나 화살촉을
% 다시 정하지 않는다 — 그것이 그림이 제각각으로 보이는 원인이다.
%
% 쓰는 법:  % gnc-style.tex — 이 볼트의 TikZ 그림이 공유하는 스타일.
%
% 저널 그림 규약을 스타일 하나로 굳혀 둔다. 그림마다 선 굵기나 화살촉을
% 다시 정하지 않는다 — 그것이 그림이 제각각으로 보이는 원인이다.
%
% 쓰는 법:  \input{gnc-style}  (figures/src/ 안에서)

\usepackage{tikz}
\usepackage{amsmath}
\usepackage{xcolor}
\usetikzlibrary{arrows.meta, positioning, calc, fit, backgrounds, decorations.pathmorphing}

% ---- 색 --------------------------------------------------------------------
% 색은 강조 하나에만 쓴다. 나머지는 검정.
\definecolor{ink}{HTML}{1A1A1A}
\definecolor{accent}{HTML}{7C3AED}
\definecolor{warn}{HTML}{B91C1C}
\definecolor{soft}{HTML}{666666}

% 기하 그림(좌표계·유도)에서만 쓰는 넷. 블록선도는 위의 흑백 + accent 로 충분하다.
\definecolor{pathc}{HTML}{0D9488}   % 계획 경로 · 항적
\definecolor{errc}{HTML}{C2410C}    % 오차 (cross-track 등)
\definecolor{alng}{HTML}{B45309}    % along-track
\definecolor{flow}{HTML}{0369A1}    % 조류 · 유속

% ---- 선과 화살촉 -----------------------------------------------------------
\tikzset{
  % 신호선. 화살촉은 작고 뾰족한 것 하나뿐이다.
  sig/.style   = {ink, line width=0.5pt, -{Stealth[length=4pt, width=3pt]}},
  wire/.style  = {ink, line width=0.5pt},
  % 강조 경로 — 파선으로 구분한다. 흑백 인쇄에서도 살아야 하기 때문이다.
  asig/.style  = {accent, line width=0.5pt, densely dashed,
                  -{Stealth[length=4pt, width=3pt]}},
  awire/.style = {accent, line width=0.5pt, densely dashed},
  % 블록. 흰 바탕, 직각, 검은 테두리.
  blk/.style   = {draw=ink, line width=0.55pt, fill=white,
                  minimum width=17mm, minimum height=8mm, inner sep=2pt, font=\small},
  % 합산점
  sum/.style   = {draw=ink, line width=0.55pt, fill=white, circle,
                  inner sep=0pt, minimum size=4.6mm, font=\scriptsize},
  asum/.style  = {draw=accent, line width=0.55pt, fill=white, circle,
                  inner sep=0pt, minimum size=4.6mm, font=\scriptsize},
  % 분기점
  dot/.style   = {ink, fill=ink, circle, inner sep=0pt, minimum size=1.5mm},
  % 신호 이름 — 선 위에 작게
  sname/.style = {font=\small, inner sep=1.5pt},
  % 그림 안의 짧은 설명. 그림 안의 글은 최소로 둔다
  note/.style  = {font=\scriptsize, soft, inner sep=1.5pt},
  anote/.style = {font=\scriptsize, accent, inner sep=1.5pt},
  % 축
  axis/.style  = {ink, line width=0.5pt},
  tick/.style  = {font=\scriptsize, soft},
}

% 캡션 한 줄 — 그림 아래 가는 선과 함께
\newcommand{\gnccaption}[2]{%
  \node[anchor=north west, text width=#1, align=left, font=\scriptsize, ink]
        at (current bounding box.south west) {\vspace{2pt}#2};}
  (figures/src/ 안에서)

\usepackage{tikz}
\usepackage{amsmath}
\usepackage{xcolor}
\usetikzlibrary{arrows.meta, positioning, calc, fit, backgrounds, decorations.pathmorphing}

% ---- 색 --------------------------------------------------------------------
% 색은 강조 하나에만 쓴다. 나머지는 검정.
\definecolor{ink}{HTML}{1A1A1A}
\definecolor{accent}{HTML}{7C3AED}
\definecolor{warn}{HTML}{B91C1C}
\definecolor{soft}{HTML}{666666}

% 기하 그림(좌표계·유도)에서만 쓰는 넷. 블록선도는 위의 흑백 + accent 로 충분하다.
\definecolor{pathc}{HTML}{0D9488}   % 계획 경로 · 항적
\definecolor{errc}{HTML}{C2410C}    % 오차 (cross-track 등)
\definecolor{alng}{HTML}{B45309}    % along-track
\definecolor{flow}{HTML}{0369A1}    % 조류 · 유속

% ---- 선과 화살촉 -----------------------------------------------------------
\tikzset{
  % 신호선. 화살촉은 작고 뾰족한 것 하나뿐이다.
  sig/.style   = {ink, line width=0.5pt, -{Stealth[length=4pt, width=3pt]}},
  wire/.style  = {ink, line width=0.5pt},
  % 강조 경로 — 파선으로 구분한다. 흑백 인쇄에서도 살아야 하기 때문이다.
  asig/.style  = {accent, line width=0.5pt, densely dashed,
                  -{Stealth[length=4pt, width=3pt]}},
  awire/.style = {accent, line width=0.5pt, densely dashed},
  % 블록. 흰 바탕, 직각, 검은 테두리.
  blk/.style   = {draw=ink, line width=0.55pt, fill=white,
                  minimum width=17mm, minimum height=8mm, inner sep=2pt, font=\small},
  % 합산점
  sum/.style   = {draw=ink, line width=0.55pt, fill=white, circle,
                  inner sep=0pt, minimum size=4.6mm, font=\scriptsize},
  asum/.style  = {draw=accent, line width=0.55pt, fill=white, circle,
                  inner sep=0pt, minimum size=4.6mm, font=\scriptsize},
  % 분기점
  dot/.style   = {ink, fill=ink, circle, inner sep=0pt, minimum size=1.5mm},
  % 신호 이름 — 선 위에 작게
  sname/.style = {font=\small, inner sep=1.5pt},
  % 그림 안의 짧은 설명. 그림 안의 글은 최소로 둔다
  note/.style  = {font=\scriptsize, soft, inner sep=1.5pt},
  anote/.style = {font=\scriptsize, accent, inner sep=1.5pt},
  % 축
  axis/.style  = {ink, line width=0.5pt},
  tick/.style  = {font=\scriptsize, soft},
}

% 캡션 한 줄 — 그림 아래 가는 선과 함께
\newcommand{\gnccaption}[2]{%
  \node[anchor=north west, text width=#1, align=left, font=\scriptsize, ink]
        at (current bounding box.south west) {\vspace{2pt}#2};}
  (figures/src/ 안에서)

\usepackage{tikz}
\usepackage{amsmath}
\usepackage{xcolor}
\usetikzlibrary{arrows.meta, positioning, calc, fit, backgrounds, decorations.pathmorphing}

% ---- 색 --------------------------------------------------------------------
% 색은 강조 하나에만 쓴다. 나머지는 검정.
\definecolor{ink}{HTML}{1A1A1A}
\definecolor{accent}{HTML}{7C3AED}
\definecolor{warn}{HTML}{B91C1C}
\definecolor{soft}{HTML}{666666}

% 기하 그림(좌표계·유도)에서만 쓰는 넷. 블록선도는 위의 흑백 + accent 로 충분하다.
\definecolor{pathc}{HTML}{0D9488}   % 계획 경로 · 항적
\definecolor{errc}{HTML}{C2410C}    % 오차 (cross-track 등)
\definecolor{alng}{HTML}{B45309}    % along-track
\definecolor{flow}{HTML}{0369A1}    % 조류 · 유속

% ---- 선과 화살촉 -----------------------------------------------------------
\tikzset{
  % 신호선. 화살촉은 작고 뾰족한 것 하나뿐이다.
  sig/.style   = {ink, line width=0.5pt, -{Stealth[length=4pt, width=3pt]}},
  wire/.style  = {ink, line width=0.5pt},
  % 강조 경로 — 파선으로 구분한다. 흑백 인쇄에서도 살아야 하기 때문이다.
  asig/.style  = {accent, line width=0.5pt, densely dashed,
                  -{Stealth[length=4pt, width=3pt]}},
  awire/.style = {accent, line width=0.5pt, densely dashed},
  % 블록. 흰 바탕, 직각, 검은 테두리.
  blk/.style   = {draw=ink, line width=0.55pt, fill=white,
                  minimum width=17mm, minimum height=8mm, inner sep=2pt, font=\small},
  % 합산점
  sum/.style   = {draw=ink, line width=0.55pt, fill=white, circle,
                  inner sep=0pt, minimum size=4.6mm, font=\scriptsize},
  asum/.style  = {draw=accent, line width=0.55pt, fill=white, circle,
                  inner sep=0pt, minimum size=4.6mm, font=\scriptsize},
  % 분기점
  dot/.style   = {ink, fill=ink, circle, inner sep=0pt, minimum size=1.5mm},
  % 신호 이름 — 선 위에 작게
  sname/.style = {font=\small, inner sep=1.5pt},
  % 그림 안의 짧은 설명. 그림 안의 글은 최소로 둔다
  note/.style  = {font=\scriptsize, soft, inner sep=1.5pt},
  anote/.style = {font=\scriptsize, accent, inner sep=1.5pt},
  % 축
  axis/.style  = {ink, line width=0.5pt},
  tick/.style  = {font=\scriptsize, soft},
}

% 캡션 한 줄 — 그림 아래 가는 선과 함께
\newcommand{\gnccaption}[2]{%
  \node[anchor=north west, text width=#1, align=left, font=\scriptsize, ink]
        at (current bounding box.south west) {\vspace{2pt}#2};}


\definecolor{knownc}{HTML}{1D4ED8}
\definecolor{ctrlc}{HTML}{B45309}

\tikzset{
  blk/.style    = {draw=ink, line width=0.9pt, fill=white, rounded corners=2pt,
                   align=center, inner sep=5pt, font=\small},
  ctrl/.style   = {blk, draw=ctrlc, line width=1.1pt, fill=alng!6},
  alloc/.style  = {blk, draw=knownc, line width=1.1pt, fill=knownc!5},
  plant/.style  = {blk, draw=pathc, line width=1.1pt, fill=pathc!10},
  outer/.style  = {draw=accent, line width=1.0pt, densely dashed, rounded corners=3pt,
                   fill=accent!4, align=center, inner sep=5pt, font=\small},
  sig/.style    = {ink, line width=0.9pt, -{Stealth[length=5.0pt,width=3.4pt]}},
  fb/.style     = {soft, line width=0.8pt, -{Stealth[length=4.6pt,width=3.2pt]}},
  bus/.style    = {soft, line width=0.8pt},
  lbl/.style    = {font=\scriptsize},
  slx/.style    = {font=\scriptsize\ttfamily, soft},
  capt/.style   = {font=\scriptsize\bfseries, warn},
  band/.style   = {font=\small\bfseries, ink},
}

\begin{document}
\begin{tikzpicture}[x=1cm, y=1cm]

% =====================================================================
% 왼쪽 — 바깥 루프. 5 주차가 채운다
% =====================================================================
\node[outer, text width=2.1cm, minimum height=4.4cm] (G) at (-2.0,0)
  {\textbf{OUTER LOOP}\\[3pt] guidance\\[1pt] \emph{Week 5}\\[6pt]
   {\scriptsize a list of waypoints becomes two commands}};

% =====================================================================
% 윗줄 — 속도 루프 (3 주차)
% =====================================================================
\draw[ink, line width=0.9pt] (0.95,1.70) circle (0.26);
\node[font=\tiny] at (0.80,1.84) {$+$};   \node[font=\tiny] at (0.95,1.52) {$-$};
\draw[sig] (-0.90,1.70) -- (0.67,1.70);
\node[lbl, anchor=south] at (-0.10,1.78) {$u_d$};

\node[ctrl, text width=4.0cm] (SP) at (3.60,1.70)
  {\textbf{SURGE SPEED LOOP} \ (PI)\\[3pt]
   {\scriptsize $X = K_{p,u}\,e_u + I$}\\[2pt]
   {\scriptsize $\dot I = K_{i,u}\,e_u + K_b\!\left(\mathrm{sat}(X)-X\right)$}\\[2pt]
   {\scriptsize \textcolor{warn}{the $K_b$ term is the anti-windup}}};
\node[slx, anchor=north, yshift=-1pt] at (SP.south) {speed loop};
\draw[sig] (1.21,1.70) -- (SP.west);
\node[lbl, anchor=south] at (1.44,1.78) {$e_u$};

% =====================================================================
% 아랫줄 — 선수각 루프 (4 주차)
% =====================================================================
\draw[ink, line width=0.9pt] (0.95,-1.70) circle (0.26);
\node[font=\tiny] at (0.80,-1.56) {$+$};  \node[font=\tiny] at (0.95,-1.88) {$-$};
\draw[sig] (-0.90,-1.70) -- (0.67,-1.70);
\node[lbl, anchor=south] at (-0.10,-1.62) {$\psi_d$};

\node[ctrl, text width=4.0cm] (HD) at (3.60,-1.70)
  {\textbf{HEADING LOOP} \ (PD)\\[3pt]
   {\scriptsize $N = K_p\,\mathrm{ssa}(e_\psi) \;-\; K_d\, r$}\\[2pt]
   {\scriptsize \textcolor{warn}{ssa wraps to $(-\pi,\pi]$}}\\[1pt]
   {\scriptsize \textcolor{knownc}{D on the measured rate}}};
\node[slx, anchor=north, yshift=-1pt] at (HD.south) {heading autopilot};
\draw[sig] (1.21,-1.70) -- (HD.west);
\node[lbl, anchor=south] at (1.44,-1.62) {$e_\psi$};

% =====================================================================
% 가운데 — 배분. 두 단계
% =====================================================================
\node[alloc, text width=4.5cm] (AL) at (9.45,0)
  {\textbf{CONTROL ALLOCATION}\\[5pt]
   {\scriptsize \emph{step 1 --- linear, invertible}}\\[2pt]
   {\scriptsize $\begin{bmatrix} F_L \\ F_R \end{bmatrix}
         = \mathbf{T}^{-1}\begin{bmatrix} X \\ N \end{bmatrix}$,\ \
         $\mathbf{T} = \begin{bmatrix} 1 & 1 \\ b & -b \end{bmatrix}$}\\[8pt]
   {\scriptsize \emph{step 2 --- nonlinear, per shaft}}\\[2pt]
   {\scriptsize $F = k\,n\lvert n\rvert \ \Rightarrow\ n = \pm\sqrt{\lvert F\rvert/k}$}};
\node[slx, anchor=north, yshift=-1pt] at (AL.south) {allocation};

\draw[sig] (SP.east) -- (6.45,1.70) -- (6.45,0.80) -- (AL.west |- 0,0.80);
\node[lbl, anchor=west] at (5.95,2.00) {$X$ \ [N]};
\draw[sig] (HD.east) -- (6.45,-1.70) -- (6.45,-0.80) -- (AL.west |- 0,-0.80);
\node[lbl, anchor=west] at (5.95,-2.00) {$N$ \ [N$\,$m]};

% =====================================================================
% 오른쪽 — 선체
% =====================================================================
\node[plant, text width=1.8cm] (OT) at (14.0,0)
  {\textbf{OTTER}\\[2pt] {\scriptsize equations of motion}};
\node[slx, anchor=north, yshift=-1pt] at (OT.south) {Otter};
\draw[sig] (AL.east) -- (OT.west);
\node[lbl, anchor=south] at (12.42,0.14) {$n_L,\ n_R$};

\draw[sig] (OT.east) -- (15.8,0);
\node[lbl, anchor=west] at (15.85,0) {$\mathbf{x}$};

% =====================================================================
% 되먹임 — 상태에서 셋을 꺼낸다
% =====================================================================
\draw[bus] (16.15,0) -- (16.6,0) -- (16.6,-4.6) -- (0.95,-4.6);
\fill[soft] (16.15,0) circle (1.4pt);
\node[lbl, soft, anchor=south east] at (16.5,-4.52)
  {$\mathbf{x} = [\,u\ \ v\ \ r\ \ N\ \ E\ \ \psi\ \cdots]^{\!\top}$};

%  u -> 속도 합산점 (아래에서) / u into the speed junction, from below
\draw[fb] (0.95,-4.6) -- (0.40,-4.6) -- (0.40,1.05) -- (0.95,1.05) -- (0.95,1.40);
\node[lbl, soft, anchor=east] at (0.33,0.40) {$u$};
%  psi -> 선수각 합산점 (아래에서) / psi into the heading junction, from below
\draw[fb] (0.95,-4.6) -- (0.95,-2.02);
\node[lbl, soft, anchor=west] at (1.02,-3.60) {$\psi$};
%  r -> 미분 항으로 직접. 오차를 미분하지 않는다 / straight into the D term
\draw[fb] (4.60,-4.6) -- (4.60,-2.92);
\fill[soft] (4.60,-4.6) circle (1.4pt);
\node[lbl, knownc, anchor=west, align=left] at (4.68,-3.80)
  {$r$ \ --- the rate is \emph{measured},\\ never $\mathrm{d}e_\psi/\mathrm{d}t$};

% =====================================================================
% 띠 설명
% =====================================================================
\node[band, anchor=west] at (-3.3,3.15) {INNER LOOP \ --- \ two controllers, one pair of propellers};
\draw[bus] (0.4,2.85) -- (11.9,2.85);
\node[capt, anchor=north, align=center] at (3.60,-5.30)
  {each loop owns one error\\ and makes one force};
\node[capt, anchor=north, align=center] at (12.0,-5.30)
  {the two forces share two propellers ---\\ that sharing is the allocation};

\end{tikzpicture}
\end{document}
